% LSPV-2 REPORT::BLOCKPROBE-TEST1::2026-10-13T15:21:43+09:00
% !TEX program = xelatex
\documentclass[11pt]{article}
\usepackage[T1]{fontenc}
\usepackage{lmodern}
\usepackage[margin=0.88in]{geometry}
\usepackage{amsmath,amssymb,amsthm,booktabs,longtable,xurl,hyperref,enumitem}
\hypersetup{hidelinks,pdftitle={BLOCKPROBE-TEST1: disjoint-block minimal-work audit}}
\newtheorem{theorem}{Theorem}
\newtheorem{lemma}[theorem]{Lemma}
\newcommand{\Ltwo}{L_{22}}
\newcommand{\OPEN}{\textbf{OPEN}}
\newcommand{\EXACT}{\textbf{PROVED/EXACT}}
\newcommand{\COMP}{\textbf{COMPUTATIONAL}}
\newcommand{\RET}{\textbf{FALSIFIED/RETIRED}}
\setlength{\emergencystretch}{3em}
\title{BLOCKPROBE-TEST1: Disjoint Compiler-Image Blocks Force Only\\Near-Linear Minimal Input Work}
\author{UNIVERSAL\_CORE\_TESTS / LSPV-2}
\date{13 October 2026}

\begin{document}
\maketitle

\begin{center}
\fbox{\parbox{0.94\linewidth}{
\textbf{Executive status.}  A same-realization byte-probe adversary yields an
unconditional $\Omega(N/\log N)$ worst-case lower bound on an explicit infinite
compiler-image subfamily of the fixed NP-complete language $\Ltwo$.  The same
subfamily has a linear-time promise recognizer.  Therefore constant-size
independent products cannot amplify this minimal-work argument to a
superpolynomial whole-computation lower bound.  The result is a candidate
theorem pending independent proof review; the bounded implementation evidence
is computational.  \textbf{P versus NP remains OPEN.}
}}
\end{center}

\section{Claim card and exact target}
The project target is the fixed language of satisfiable encodings of nonempty,
proper, distinct-clause signed exact-$(2,2)$ 3-CNFs; malformed strings are
rejected.  The desired endpoint is
\[
 \forall A\,\forall k\in\mathbb N\,\forall C>0\,\forall N_0\,
 \exists G\in\mathcal C_{22}:\quad
 |\operatorname{enc}G|\ge N_0\ \land\
 T_A(G)>C|\operatorname{enc}G|^k .
\]
This report does not prove that endpoint.  It tests the narrower candidate
that many independent compiler-image blocks force superpolynomial minimal
work for every deterministic decider.

\section{Construction}
Let $Y$ be the six-vertex graph with edges
\[
 02,03,05,12,13,14,23
\]
and let $Z$ have edges
\[
 01,02,03,12,13,23,45.
\]
Both have labelled degree vector $(3,3,3,3,1,1)$ and seven edges.  Exhaustive
enumeration gives 24 proper 3-colourings of $Y$ and none of $Z$.
Apply the actual repository pipeline
\[
G\mapsto \operatorname{colouring\_cnf}(G)
\mapsto \operatorname{properize23}(G)
\mapsto U(G)=\operatorname{compile22}(G).
\]
The proper sources for $Y,Z$ each have 57 variables and 84 clauses.  For
$r\ge1$, concatenate $r$ variable-disjoint copies before applying $U$.
Write $X_r$ for the all-$Y$ target and $X_{r,i}$ for the target with block $i$
replaced by $Z$.  Then $X_r\in\Ltwo$ and $X_{r,i}\notin\Ltwo$.

\section{Candidate theorem: exact byte-probe baseline}
Fix the compact ordered JSON byte encoding used by the audit.  Let $D_i$ be
the byte positions on which $\operatorname{enc}X_r$ and
$\operatorname{enc}X_{r,i}$ differ.

\begin{theorem}[CANDIDATE pending independent review]\label{thm:blockprobe}
For every $r$, the sets $D_1,\ldots,D_r$ are nonempty and pairwise disjoint.
Every deterministic adaptive byte-query algorithm that is correct on all
valid $\Ltwo$ encodings queries at least $r$ bytes on
$\operatorname{enc}X_r$.  Under this encoding,
\[
  N_r=|\operatorname{enc}X_r|=\Theta(r\log r),
\]
so this gives $\Omega(N_r/\log N_r)$ byte probes and hence the same lower
bound on any model charging unit cost per byte probe.
\end{theorem}

\begin{proof}
The compiler preserves the proper source clauses as an initial ordered prefix.
Because the two block templates have the same source variable and clause
counts and the same labelled degree profile, the balancing deficits, clone
counts, gadget layout, total variable count and total clause count agree.
Changing block $i$ changes a nonempty set of bytes inside that block's
source-prefix interval and no other block's interval.  Hence the $D_i$ are
nonempty and disjoint.

Run a deterministic byte-query algorithm on $X_r$.  If it never queries a
position in $D_i$, every answer in its adaptive transcript is identical on
$X_r$ and $X_{r,i}$.  It must then return the same bit on a YES and a NO input,
a contradiction.  Thus the run on the single realization $X_r$ intersects
every $D_i$ and makes at least $r$ probes.  Target variables and clauses are
linear in $r$; printed variable identifiers need $\Theta(\log r)$ digits, so
the selected byte encoding has length $\Theta(r\log r)$.
\end{proof}

\paragraph{Model boundary.}
This is a byte-query theorem, not an unrestricted RAM, circuit, arithmetic,
proof-system, or time lower bound.  Wider word reads require a separately
stated word-size accounting.  Query complexity and full computation time are
not silently identified.

\section{The matching algorithmic escape}
The promise support $\{X_r,X_{r,i}:1\le i\le r\}$ has a direct
$O(N_r)$ recognizer: scan the source-clause prefix block by block and compare
each block with the fixed $Y/Z$ templates.  This charges parsing, access and
comparison.  It does not decide arbitrary $\Ltwo$ inputs, but it proves that
the direct product does not create a hard promise.  OR and parity give the
same structural warning: block sensitivity can be linear while a single scan
is polynomial.

Therefore the proposed amplification
\[
 \text{polynomially many independent constant-size hard-looking blocks}
 \Longrightarrow \text{superpolynomial whole-machine time}
\]
is \RET.  The positive input-probe baseline in
Theorem~\ref{thm:blockprobe} remains intact.

\section{Bounded reproducible audit}
The script
\path{experiments/review/l22_disjoint_block_probe_audit.py}
ran for 6.18 seconds under a 30-second shell timeout and a cooperative
25-second deadline.  No chunking was needed.  Its SHA-256 is
\texttt{f4c7342039e11966547f421f8a4b1addcaa13e0f3ae7491dcdc4ecb766a056b9}.
The raw result is
\path{results/raw/review/2026-10-13_1521_l22_disjoint_block_probe.json},
SHA-256
\texttt{fc00a7c0d85e5d85e65e63418dccc955397039deeeb4365c63449095705fff00},
semantic digest
\texttt{48828f0c71788cd6189c63dd38bfe62d84ac9b74896578bc8907559d474f2919}.

\begin{center}
\begin{tabular}{rrrrr}
\toprule
$r$ & encoded bytes & target variables & target clauses & disjoint hits\\
\midrule
1 & 367345 & 14364 & 19152 & 1\\
2 & 779113 & 28728 & 38304 & 2\\
3 & 1190881 & 43092 & 57456 & 3\\
4 & 1602649 & 57456 & 76608 & 4\\
\bottomrule
\end{tabular}
\end{center}

All ten single-bad-block variants were valid exact-$(2,2)$ targets with the
same target counts and byte length as their all-YES mate.  The audit found
2475 differing bytes and passed all ten pairwise-disjointness checks.  Twelve
COPY, interpreter/re-encode and DEAD-suffix wrapper controls passed.  It also
ran concrete small controls for 2SAT, Horn forward semantics, affine Gaussian
elimination, bounded-treewidth path DP, independent products, redundant
clauses, tautological padding, compact/pretty encodings, OR and parity.  These
are finite controls only; no asymptotic hardness is inferred.

\section{SAME\_REALIZATION and cost audit}
The lower-bound transcript is the actual run on $X_r$; the $X_{r,i}$ are only
adversary witnesses.  Construction and discovery are not credited as work
performed by the tested decider.  The explicit promise recognizer is charged
for decoding, reading every source-prefix byte, template comparisons and its
one-bit output.  Reconstruction is unnecessary because the decision language
asks only for a bit.  This removes any appeal to mandatory compatibility
extraction, proof emission or hidden semantic state.

\section{Literature and barrier-transfer preflight}
\begin{itemize}[leftmargin=*]
\item \textbf{Directly transferable theorem/model.}  Nisan's block-sensitivity
work relates block sensitivity and deterministic decision-tree complexity in
an explicit query model.  The elementary transcript proof above is stated
directly for our byte-query model and does not import a whole-machine bound.
\item \textbf{Useful analogy requiring a new transfer theorem.}  Rivest and
Vuillemin prove substantial adjacency-query lower bounds for monotone graph
properties.  Our property, representation and machine model differ; their
result motivates the probe audit but does not lower-bound arbitrary $\Ltwo$
computation.
\item \textbf{Known theorem with incompatible resource promise.}  Fortnow's
SAT time--space tradeoffs and later refinements give unconditional bounds only
under joint small-space/time hypotheses.  They do not cover an arbitrary
polynomial-time decider with unrestricted polynomial space.
\item \textbf{Barrier.}  Aaronson--Wigderson algebrization reinforces that an
algebraic or query invariant is not by itself a general Boolean lower bound.
No nonrelativizing transfer is supplied here.
\end{itemize}

\section{Post-quantum technique ledger}
Schur squares, quadratic relation ranks and Hilbert/syzygy invariants are not
used: the candidate is defeated before algebraic structure becomes relevant.
Collision entropy is an \textbf{ANALOGY/CANDIDATE}: independent blocks give a
transparent collision partition, but its effective entropy is exploitable by
the linear recognizer.  LWE-style worst-case/average-case architecture is a
\textbf{FAILED TRANSFER} here because no balanced hard distribution or charged
reduction is present.  The same visible structure is an algorithmic weakness.

\section{CORPUS REVIEWED}
The run surveyed the current repository tree and read root \path{main.tex},
\path{AGENTS.md}, \path{CODEX.md}, the complete
\path{.codex/skills/pnp-research/SKILL.md} and referenced protocol files, all
current task specifications, actual \path{main4.tex}/\path{main5.tex}, newest
verified MAIN \path{research/main/2026-10-13_1444_MAIN_RS5REVIEW1.tex}, the
newest LSPV-1--4 reports, SINGLE22 and EARLY22 governance notes, the original
\path{research/agents/lspv-4/2026-10-13_0035_LSPV-4_PC22_ALLFIELD1.tex},
\path{universal22.py}, \path{audit_universal22.py}, relevant historical
compiler/projection reports, and the HEC/VSC-HEC/G-PPC-HEC/CRS-PPC/CCSPC
arbitrary-algorithm branch.  Repository searches found no prior
BLOCKPROBE formulation.

\section{DUPLICATION / CONTRADICTION AUDIT}
This result does not alter COMPOSEPROJ1, the exact compiler, the exact-(2,2)
minimality theorem, Resolution projection, or PC/PCR projection.  It is
consistent with NOBIASFW1, EQSPECTRAL1 and COSPECBLOWRED1: conspicuous
structure may make a promised compiler-image family easy.  It contradicts
only an overbroad inference from many independent sensitive blocks to
superpolynomial computation.  Expanded-monomial proof size is still not
arbitrary runtime.

\section{ADOPTED / REJECTED / RETIRED PRIOR COMPONENTS}
\begin{itemize}[leftmargin=*]
\item \textbf{ADOPTED:} the verified graph-colouring adapter, transparent
seven-clause padding, exact-$(2,2)$ compiler, explicit encoding and invalid
input rejection convention.
\item \textbf{ADOPTED:} the requirement that a direct proof for the one fixed
language $\Ltwo$ needs no later broad-SAT transfer.
\item \textbf{REJECTED:} mandatory extraction of $H(F)$, a PC proof, a DPLL
certificate, or a universal implementation compiler.
\item \textbf{RETIRED:} constant-size independent-block amplification as a
route from input sensitivity to superpolynomial whole-machine time.
\end{itemize}

\section{EARLY22-INTEGRATION1 acknowledgment}
Sources actually read include the verified governance memo, original compiler
and projection reports surfaced by its ledger, current supersession notices,
and the graph-colouring successor construction.  The exact-(2,2) compiler,
minimality and restricted Resolution results are integrated only within their
hypotheses; the monotone cubic-pair result and affine-continuation limitation
remain separate.  Exact location: this section and the status/root notes below.
Next check: test whether a growing-block replacement preserves disjoint
sensitivity after canonical decoding while defeating template recognition.

\section{SINGLE22-SEPARATION1 acknowledgment}
The tested source is the graph-colouring/transparent-seven-clause report at
the full path and blob specified by governance, plus current repository
\path{universal22.py}; the earlier selector/cardinality body was not
substituted.  The fixed language remains $\Ltwo$, not one formula or a NO-only
benchmark.  What changed is exact: we now have a direct all-decider
byte-probe baseline on an infinite valid family, and a matching polynomial
promise algorithm that sharply limits its use.  The whole-computation arrow
needed for $\Ltwo\notin\mathsf P$ remains open.

\section{WHAT CHANGED RELATIVE TO ROOT MAIN}
The current default-branch \path{main.tex} readback is a markerless placeholder
with blob SHA \texttt{725bf15b399d4efbc3659164bda03f0ebda2cb16}.
Because the required LSPV-2 ownership markers are absent, AGENTS.md forbids a
specialist from restructuring it.  No root write is attempted.  LSPV-5 must
restore the living document from a verified MAIN snapshot before specialist
reconciliation can resume.

\section{STALE OR OVERSTATED CLAIMS FOUND}
The root placeholder is stale relative to verified MAIN and specialist
artifacts.  Any statement that block sensitivity, adjacency probes, number of
semantic states, or explicit-monomial proof size already lower-bounds every
deterministic $\Ltwo$ computation is overstated and must remain retired.

\section{Predecessor write recovery}
The LSPV-1 artifact
\path{research/agents/lspv-1/2026-10-13_1409_LSPV-1_TRACEAXIOM-FIREWALL1.tex}
verifies.  No canonical artifact exists for the immediately following LSPV-1
cycle.  Its exact unfinished body is not available from GitHub, so this run
issues a recovery request to preserve and retry that body without reconstructing
scientific content.  Failure classification: \textbf{UNKNOWN / unpublished
body}; the current run is not blocked.

\section{Status and next obligation}
\begin{center}
\begin{tabular}{p{0.20\linewidth}p{0.72\linewidth}}
\toprule
Status & Result\\
\midrule
CANDIDATE & Theorem~\ref{thm:blockprobe}, pending independent proof review.\\
COMPUTATIONAL & Exact compiler/encoding checks for $r\le4$ and hostile controls.\\
FALSIFIED/RETIRED & Constant-size independent blocks as superpolynomial amplification.\\
OPEN & A growing internally hard and nonlocally coupled valid family, plus a
model-matched whole-computation lower bound for every total deterministic
uniform decider and every polynomial exponent.\\
\bottomrule
\end{tabular}
\end{center}

The highest-value next test is a growing-block family with equalized syntax
whose decoded quotient is not fixed, followed by canonicalization, WL,
Gaussian, flow/counting, bounded-treewidth and template-recognition attacks
before compilation.  \textbf{P versus NP remains OPEN.}

\begin{thebibliography}{9}
\bibitem{Nisan1991}
N.~Nisan, ``CREW PRAMs and Decision Trees,'' \emph{SIAM Journal on Computing}
20(6):999--1007, 1991. DOI: \url{https://doi.org/10.1137/0220062}.
\bibitem{RV1976}
R.~L. Rivest and J.~Vuillemin, ``On recognizing graph properties from
adjacency matrices,'' \emph{Theoretical Computer Science} 3(3):371--384,
1976. DOI: \url{https://doi.org/10.1016/0304-3975(76)90053-0}.
\bibitem{Fortnow2000}
L.~Fortnow, ``Time--Space Tradeoffs for Satisfiability,'' \emph{Journal of
Computer and System Sciences} 60(2):337--353, 2000.
DOI: \url{https://doi.org/10.1006/jcss.1999.1671}.
\bibitem{AW2009}
S.~Aaronson and A.~Wigderson, ``Algebrization: A New Barrier in Complexity
Theory,'' \emph{ACM Transactions on Computation Theory} 1(1), 2009.
DOI: \url{https://doi.org/10.1145/1490270.1490272}.
\end{thebibliography}

\end{document}
% END_OF_VERIFIED_REPORT::LSPV-2::BLOCKPROBE-TEST1
